\chapter{Estimation de la covariance et matrices aléatoires}\label{ch:qm:covariance-estimation-and-random-matrices}

Un \omterm{def:qm:covariance-estimation-and-random-matrices:pca}{portefeuille de variance minimale} de 200 actions, construit à partir de la \omterm{def:qm:covariance-estimation-and-random-matrices:cov}{matrice de covariance empirique} de deux années de rendements quotidiens, promet une volatilité annuelle de
5.7~\%. Conservé l'année suivante, il affiche 9.6~\%. Rien n'a cédé, et le marché n'a pas changé~: dans la simulation du chapitre, les rendements des deux années
proviennent de la même covariance factorielle, dont le vrai \omterm{def:qm:covariance-estimation-and-random-matrices:pca}{portefeuille de variance minimale} présente une volatilité de 7.4~\%. Avec 200 actifs et 500 jours, la
\omterm{def:qm:covariance-estimation-and-random-matrices:cov}{matrice de covariance empirique} compte 20\,100 paramètres estimés à partir de 100\,000 nombres, et l'optimiseur trouve les directions où ses erreurs font paraître le risque
le plus faible. Ce chapitre porte sur ce qu'est une \omterm{def:qm:covariance-estimation-and-random-matrices:cov}{matrice de covariance empirique} en grande dimension~: la loi que suivent ses valeurs propres lorsqu'il n'y a que du bruit,
les quelques valeurs propres qui portent l'information, et les \omterm{def:qm:estimation:estimator}{estimateurs} qui nettoient le reste (rétrécissement vers une cible, valeur propre par valeur propre, écrêtage,
et structure factorielle), jugés à l'aune du seul test qui compte pour un portefeuille~: le risque réalisé face au risque promis.

\section{La matrice de covariance empirique en grande dimension}

\begin{definition}[Matrice de covariance, matrice de covariance empirique]\label{def:qm:covariance-estimation-and-random-matrices:cov}
La \emph{matrice de covariance}\index{matrice de covariance} d'un vecteur aléatoire $r \in \R^N$ est $\Sigma = \E[(r - \mu)(r - \mu)^\top]$. À partir de $T$ observations $r_1, \dots, r_T$
de moyenne empirique $\bar r$, la \emph{matrice de covariance empirique}\index{matrice de covariance empirique} est $S = \frac1T\sum_t(r_t - \bar r)(r_t - \bar r)^\top$ (ou avec le diviseur
$T - 1$).
\end{definition}

\begin{definition}[Analyse en composantes principales, portefeuille de variance minimale]\label{def:qm:covariance-estimation-and-random-matrices:pca}
\emph{analyse en composantes principales}\index{analyse en composantes principales} (ACP) diagonalise une \omterm{def:qm:covariance-estimation-and-random-matrices:cov}{matrice de covariance} (ou de corrélation), $S = U\Lambda U^\top$~: les vecteurs propres sont
des portefeuilles non corrélés, et les valeurs propres leurs variances, par ordre décroissant. Le \emph{portefeuille de variance minimale}\index{portefeuille de variance minimale} pour une
covariance $\Sigma$ est $w = \Sigma^{-1}\mathbf 1/(\mathbf 1^\top\Sigma^{-1}\mathbf 1)$, le portefeuille entièrement investi de plus petite variance, $1/(\mathbf 1^\top\Sigma^{-1}\mathbf 1)$.
\end{definition}

Le \omterm{def:qm:covariance-estimation-and-random-matrices:pca}{portefeuille de variance minimale} n'a pas besoin de rendements espérés, ce qui en fait le test le plus net d'une estimation de covariance~; il inverse aussi la matrice, ce qui en fait le plus
exigeant. $\Sigma^{-1}$ pondère chaque vecteur propre par l'inverse de sa valeur propre, si bien que les plus petites valeurs propres empiriques, celles que le bruit déprime
le plus, reçoivent le poids le plus fort. Avec $q = N/T$ fixé lorsque les deux croissent, et des rendements gaussiens, le portefeuille construit à partir de $S$ a une variance dans l'échantillon d'environ $(1 - q)$
fois le minimum vrai, et une variance vraie d'environ $1/(1 - q)$ fois ce minimum~: le rapport entre volatilité réalisée et volatilité prédite vaut environ $1/(1 - q)$.

Le modèle vrai du chapitre comporte 200 actions, un facteur de marché et cinq secteurs de quarante actions, avec des volatilités spécifiques comprises entre 17~\% et 44~\% par an. Sur cinquante
historiques simulés sur deux ans ($q = 0.4$), le \omterm{def:qm:covariance-estimation-and-random-matrices:pca}{portefeuille de variance minimale} empirique annonce 5.71~\% et délivre 9.58~\% face à la valeur vraie (9.61~\% mesurés sur les 250
jours suivants)~: un rapport de 1.68 face à $1/(1 - 0.4) = 1.67$ pour la formule. En faisant varier l'historique de huit ans à moins d'un an ($q$ de 0.1 à 0.9), le rapport simulé suit
la formule sur toute la plage, atteignant 10.1 pour $q = 0.9$ (la \cref{fig:qm:covariance-estimation-and-random-matrices:bias}).

\begin{omfigure}
\begin{tikzpicture}
\begin{semilogyaxis}[width=10cm,height=5cm,xlabel={$q = N/T$ (200 actions)},ylabel={volatilité vraie / prédite},
  tick label style={font=\small},label style={font=\small},xmin=0,xmax=0.95,ymin=1,ymax=12,grid=major,grid style={gray!25},
  ytick={1,1.5,2,3,5,10},yticklabels={1,1.5,2,3,5,10},
  legend style={font=\footnotesize,at={(0.03,0.97)},anchor=north west,fill=white,draw=none},table/col sep=comma]
\addplot[only marks,mark=*,mark size=2pt,omDef] table[x=q,y=sim]{figdata/methods/22-covariance-estimation-and-random-matrices/bias.csv};
\addplot[omThm,thick,domain=0:0.93,samples=60] {1/(1-x)};
\legend{simulé (20 historiques chacun),$1/(1 - q)$}
\end{semilogyaxis}
\end{tikzpicture}

\omcaption{Rapport entre la volatilité vraie et la volatilité prédite du \omterm{def:qm:covariance-estimation-and-random-matrices:pca}{portefeuille de variance minimale} empirique de 200 actions, en fonction de $q = N/T$ (historiques de 2\,000 à 222 jours),
et la formule en grande dimension $1/(1 - q)$. Données~: le tutoriel du chapitre, graine fixée.}\label{fig:qm:covariance-estimation-and-random-matrices:bias}
\end{omfigure}

\section{La loi des valeurs propres du bruit pur}

\begin{theorem}[Loi de Marchenko--Pastur]\label{thm:qm:covariance-estimation-and-random-matrices:mp}
Si la matrice $T \times N$ des rendements a des entrées i.i.d. de moyenne nulle et de variance $\sigma^2$ et que $N/T \to q \in (0, 1)$, la distribution empirique des valeurs propres de
la \omterm{def:qm:covariance-estimation-and-random-matrices:cov}{matrice de covariance empirique} converge vers la \emph{loi de Marchenko--Pastur}\index{loi de Marchenko--Pastur}, de densité
\[
\rho(\lambda) = \frac{\sqrt{(\lambda_+ - \lambda)(\lambda - \lambda_-)}}{2\pi q\sigma^2\lambda}, \qquad \lambda_\pm = \sigma^2(1 \pm \sqrt q)^2,
\]
sur $[\lambda_-, \lambda_+]$ (Marchenko et Pastur, 1967).
\end{theorem}

Les valeurs propres vraies valent toutes $\sigma^2$~; l'échantillon les disperse de $0.14\sigma^2$ à $2.66\sigma^2$ pour $q = 0.4$. Les plus petites sont trop petites d'un facteur sept, et ce sont
les directions que l'optimiseur de variance minimale affectionne. Sur 500 jours de 200 séries gaussiennes centrées réduites indépendantes, les valeurs propres de la matrice de corrélation empirique vont de 0.144 à 2.565,
à l'intérieur des bornes théoriques 0.135 et 2.665 (la \cref{fig:qm:covariance-estimation-and-random-matrices:mp})~; Laloux, Cizeau, Bouchaud et Potters (1999) ont trouvé que la plupart des
valeurs propres des matrices de corrélation réelles d'actions se situent dans la même loi.

\begin{omfigure}
\begin{tikzpicture}
\begin{axis}[width=10cm,height=5cm,xlabel={valeur propre de la matrice de corrélation empirique},ylabel={densité},
  tick label style={font=\small},label style={font=\small},xmin=0,xmax=3,ymin=0,ymax=1.2,grid=major,grid style={gray!25},
  legend style={font=\footnotesize,at={(0.97,0.97)},anchor=north east,fill=white,draw=none},table/col sep=comma]
\addplot[omDef,thick,const plot mark mid,fill=omDef!15] table[x=x,y=noise]{figdata/methods/22-covariance-estimation-and-random-matrices/mp_hist.csv};
\addplot[black,very thick] table[x=x,y=density]{figdata/methods/22-covariance-estimation-and-random-matrices/mp_curve.csv};
\addplot[omThm,thick,const plot mark mid] table[x=x,y=factor]{figdata/methods/22-covariance-estimation-and-random-matrices/mp_hist.csv};
\legend{bruit pur,Marchenko--Pastur ($q = 0.4$),modèle factoriel (cœur du spectre)}
\end{axis}
\end{tikzpicture}

\omcaption{Valeurs propres de matrices de corrélation empiriques de 200 séries sur 500 jours~: bruit indépendant (en grisé) face à la densité de Marchenko--Pastur, et celles du \omterm{def:qm:covariance-estimation-and-random-matrices:factor}{modèle factoriel}
inférieures à 3 (sa valeur propre de marché, 51.4, et trois de ses valeurs propres sectorielles sortent du graphique). Données~: le tutoriel du chapitre, graine fixée.}\label{fig:qm:covariance-estimation-and-random-matrices:mp}
\end{omfigure}

\begin{definition}[Modèle de covariance à pics]\label{def:qm:covariance-estimation-and-random-matrices:spiked}
Dans le \emph{modèle de covariance à pics}\index{modèle de covariance à pics} (Johnstone, 2001), la covariance de population est l'identité, sauf pour quelques valeurs propres $\ell_1, \dots, \ell_k
> 1$, les pics, qui représentent des facteurs au-dessus d'un plancher de bruit.
\end{definition}

Un pic n'est visible que s'il est assez grand~: lorsque $\ell \le 1 + \sqrt q$, la plus grande valeur propre empirique reste collée au bord $(1 + \sqrt q)^2$~; au-delà de ce seuil, elle se détache et
converge vers $\ell(1 + q/(\ell - 1))$, supérieure à $\ell$ (Baik, Ben Arous et Péché, 2005, pour les données complexes~; Baik et Silverstein, 2006, pour les données réelles). Pour $q = 0.4$, le seuil est
1.63~: un facteur expliquant moins de 1.63 fois la variance moyenne d'une action est invisible sur deux ans de données quotidiennes, aussi réel soit-il,
(la \cref{fig:qm:covariance-estimation-and-random-matrices:spike}). Dans le \omterm{def:qm:covariance-estimation-and-random-matrices:factor}{modèle factoriel}, cinq valeurs propres empiriques dépassent le bord~: celle du marché, 51.4 (population 52.9) et quatre
valeurs propres sectorielles de 2.96 à 3.68 (population de 2.83 à 2.99), chacune poussée vers le haut comme l'indique la formule.

\begin{omfigure}
\begin{tikzpicture}
\begin{axis}[width=10cm,height=5cm,xlabel={pic de population $\ell$},ylabel={plus grande valeur propre empirique},
  tick label style={font=\small},label style={font=\small},xmin=1,xmax=5.2,ymin=2,ymax=6,grid=major,grid style={gray!25},
  legend style={font=\footnotesize,at={(0.03,0.97)},anchor=north west,fill=white,draw=none},table/col sep=comma]
\addplot[only marks,mark=*,mark size=2pt,omDef] table[x=ell,y=sim]{figdata/methods/22-covariance-estimation-and-random-matrices/spike.csv};
\addplot[omThm,thick,domain=1.633:5.2,samples=60] {x*(1+0.4/(x-1))};
\addplot[omThm,thick,domain=1:1.633,samples=2] {2.665};
\addplot[black,dashed,thin,domain=1:5.2,samples=2] {x};
\legend{{simulé ($N = 200$, $T = 500$)},théorie,$y = \ell$}
\end{axis}
\end{tikzpicture}

\omcaption{La transition du pic~: la plus grande valeur propre empirique en fonction d'un unique pic de population $\ell$ (toutes les autres valeurs propres valant un), avec la théorie~: le bord du spectre $(1 + \sqrt q)^2 =
2.66$ sous le seuil $1 + \sqrt q = 1.63$, et $\ell(1 + q/(\ell - 1))$ au-dessus. Données~: le tutoriel du chapitre, graine fixée.}\label{fig:qm:covariance-estimation-and-random-matrices:spike}
\end{omfigure}

\section{Rétrécissement linéaire et non linéaire}

\begin{definition}[Rétrécissement linéaire]\label{def:qm:covariance-estimation-and-random-matrices:linear}
Le \emph{rétrécissement linéaire}\index{rétrécissement linéaire} (linear shrinkage) remplace $S$ par $\delta F + (1 - \delta)S$ pour une cible structurée $F$, comme un multiple de l'identité ou la matrice à corrélation
constante, avec une intensité $\delta \in [0, 1]$ choisie pour minimiser la distance de Frobenius espérée à $\Sigma$.
\end{definition}

C'est l'idée de James--Stein du chapitre~14 appliquée à une matrice. Ledoit et Wolf (2004) ont dérivé des estimations convergentes de l'intensité optimale, qui ne nécessite que les données. Pour
la cible $mI$, $m = \operatorname{tr}S/N$, l'intensité vaut $\min(\bar b^2, d^2)/d^2$, avec $d^2 = \lVert S - mI\rVert^2/N$ et $\bar b^2 = T^{-2}\sum_t\lVert r_tr_t^\top - S\rVert^2/N$, la dispersion de $S$ autour de la
cible, rapportée à son propre bruit d'échantillonnage. Sur les historiques du chapitre, l'intensité vers l'identité vaut en moyenne 0.034 et fait passer la volatilité réalisée de 9.61~\% à seulement
9.14~\%~: l'identité est une cible médiocre lorsque les actions partagent un facteur de marché. Vers la cible à corrélation constante, l'intensité vaut 0.178, la prévision devient honnête (7.64~\%), mais la
volatilité réalisée reste à 9.14~\%.

\begin{definition}[Estimateur équivariant par rotation, rétrécissement non linéaire]\label{def:qm:covariance-estimation-and-random-matrices:nonlinear}
Un \emph{estimateur équivariant par rotation}\index{estimateur équivariant par rotation} (rotation-equivariant estimator) conserve les vecteurs propres empiriques et ne modifie que les valeurs propres~: $\hat\Sigma = U\operatorname{diag}(d)U^\top$.
Le \emph{rétrécissement non linéaire}\index{rétrécissement non linéaire} (nonlinear shrinkage) choisit chaque $d_i$ séparément, comme estimation de la valeur oracle $u_i^\top\Sigma u_i$, la variance vraie du $i$-ième
portefeuille propre empirique.
\end{definition}

Les valeurs propres oracle sont le meilleur résultat qu'un \omterm{def:qm:covariance-estimation-and-random-matrices:nonlinear}{estimateur équivariant par rotation} puisse atteindre (elles minimisent l'erreur de Frobenius, à $U$ donné). Elles sont inconnues, mais la théorie des matrices aléatoires les exprime
à travers la transformée de Stieltjes de la densité limite des valeurs propres, que l'on peut estimer à partir des valeurs propres empiriques elles-mêmes~; Ledoit et Wolf (2020) ont donné une
version explicite à l'aide d'un \omterm{def:qm:estimation:estimator}{estimateur} de densité à noyau et de sa transformée de Hilbert. Sur un test dont les valeurs propres de population s'étalent de 0.5 à 2, elle réduit l'erreur quadratique moyenne des valeurs propres
par rapport à l'oracle, de 0.49 à 0.0034. Sur le test de portefeuille, elle prévoit 7.36~\% et réalise 8.25~\%.

\section{Écrêtage des valeurs propres et structure factorielle}

\begin{definition}[Écrêtage des valeurs propres]\label{def:qm:covariance-estimation-and-random-matrices:clip}
L'\emph{écrêtage des valeurs propres}\index{ecretage des valeurs propres@écrêtage des valeurs propres} (eigenvalue clipping~; Laloux et al., 1999) conserve les valeurs propres de la matrice de corrélation empirique situées au-dessus du bord du spectre de Marchenko--Pastur et remplace toutes
les autres par leur moyenne, ce qui préserve la trace, puis restaure les variances empiriques.
\end{definition}

\begin{definition}[Modèle factoriel]\label{def:qm:covariance-estimation-and-random-matrices:factor}
Un \emph{modèle factoriel}\index{modèle factoriel} (factor model) de covariance est $\Sigma = BF B^\top + D$, avec $k \ll N$ expositions aux facteurs $B$, une covariance des facteurs $F$ et une variance spécifique diagonale $D$~; les modèles
factoriels statistiques prennent $B$ dans les premières composantes principales de $S$.
\end{definition}

L'écrêtage prévoit 7.27~\% et réalise 8.05~\%~; un modèle à six facteurs en composantes principales, avec un résidu diagonal, prévoit 6.65~\% et réalise 7.74~\%, le plus proche de l'optimum vrai de 7.44~\%
(la \cref{fig:qm:covariance-estimation-and-random-matrices:estimators}). Les rapports entre volatilité réalisée et volatilité prévue valent 1.68 (échantillon), 1.51 (linéaire vers l'identité), 1.19 (linéaire vers
la corrélation constante), 1.12 (non linéaire), 1.11 (écrêtage) et 1.17 (\omterm{def:qm:covariance-estimation-and-random-matrices:factor}{modèle factoriel}). Le \omterm{def:qm:covariance-estimation-and-random-matrices:factor}{modèle factoriel} l'emporte ici, car la vraie structure est un \omterm{def:qm:covariance-estimation-and-random-matrices:factor}{modèle factoriel} à six facteurs, que les composantes principales
retrouvent~; avec un modèle vrai moins net, l'\omterm{def:qm:estimation:estimator}{estimateur} non linéaire, qui ne fait aucune hypothèse de structure, reste le choix par défaut le plus prudent. Fan, Liao et Mincheva (2013) combinent les deux~: des facteurs
pour les pics, seuillage pour le reste.

\begin{omfigure}
\begin{tikzpicture}
\begin{axis}[width=12.5cm,height=5cm,ybar,bar width=9pt,
  xtick={0,1,2,3,4,5},xticklabels={échantillon,{LW\\identité},{LW corr.\\constante},non linéaire,écrêtage,{ACP à\\6 facteurs}},
  xticklabel style={font=\footnotesize,align=center},
  ylabel={volatilité annuelle (\%)},tick label style={font=\small},label style={font=\small},
  xmin=-0.6,xmax=5.6,ymin=0,ymax=11,grid=major,grid style={gray!25},
  legend columns=3,legend style={font=\footnotesize,at={(0.5,-0.3)},anchor=north,draw=none,fill=none,
  /tikz/every even column/.append style={column sep=8pt}},table/col sep=comma]
\addplot[fill=omDef!55,draw=omDef] table[x=k,y=pred]{figdata/methods/22-covariance-estimation-and-random-matrices/estimators.csv};
\addplot[fill=omThm!50,draw=omThm] table[x=k,y=real]{figdata/methods/22-covariance-estimation-and-random-matrices/estimators.csv};
\addplot[black,dashed,thick,sharp plot,forget plot] coordinates {(-0.6,7.44) (5.6,7.44)};
\addlegendimage{line legend,black,dashed,thick}
\legend{prévue (en échantillon),réalisée (année suivante),optimum vrai}
\end{axis}
\end{tikzpicture}

\omcaption{\omterm{def:qm:covariance-estimation-and-random-matrices:pca}{Portefeuilles de variance minimale} de 200 actions construits à partir de six \omterm{def:qm:estimation:estimator}{estimateurs} de covariance sur deux ans de rendements quotidiens~: la volatilité que chacun prévoit et la volatilité que chacun réalise sur
l'année suivante, moyennées sur cinquante historiques simulés~; la ligne en tirets représente l'optimum vrai, 7.44~\%. Données~: le tutoriel du chapitre, graine fixée.}\label{fig:qm:covariance-estimation-and-random-matrices:estimators}
\end{omfigure}

\section{Tutoriel~: le portefeuille qui promettait 6~\%}\label{tut:qm:covariance-estimation-and-random-matrices:promise}

\begin{tutorial}
\textbf{Objectif.} Construire des \omterm{def:qm:covariance-estimation-and-random-matrices:pca}{portefeuilles de variance minimale} à partir de six \omterm{def:qm:estimation:estimator}{estimateurs} de covariance et comparer ce qu'ils promettent à ce qu'ils réalisent.
\textbf{État final~:} les \cref{fig:qm:covariance-estimation-and-random-matrices:bias,fig:qm:covariance-estimation-and-random-matrices:estimators}~; le ratio de 1.68 pour le portefeuille empirique et
celui des meilleurs \omterm{def:qm:estimation:estimator}{estimateurs}, de 1.11 à 1.19.
\begin{tutsteps}
\item \textbf{\omterm{def:qm:covariance-estimation-and-random-matrices:nonlinear}{Rétrécissement non linéaire} analytique}~: densité à noyau des valeurs propres, sa transformée de Hilbert et les valeurs propres rétrécies.
\omcode{code/firm/covest/firm_covest.py}{86}{109}{Rétrécissement non linéaire analytique des valeurs propres empiriques.}\label{lst:qm:covariance-estimation-and-random-matrices:nl}
\item \textbf{Écrêtage} au bord du spectre de Marchenko--Pastur.
\omcode{code/firm/covest/firm_covest.py}{112}{134}{Écrêtage des valeurs propres de la matrice de corrélation empirique.}\label{lst:qm:covariance-estimation-and-random-matrices:clip}
\item \textbf{Exécuter} \texttt{evaluate()}, \texttt{bias\_vs\_q()}, \texttt{spectrum()} et \texttt{spike()} dans \texttt{qm\_covest.py}, puis \texttt{fig\_covest.py}.
\end{tutsteps}
\textbf{Pour aller plus loin.} Ajouter une contrainte long only et observer quelle part de l'erreur du portefeuille empirique elle supprime (elle agit comme un rétrécissement)~; remplacer les rendements gaussiens par des
rendements de Student-$t$ et comparer de nouveau les \omterm{def:qm:estimation:estimator}{estimateurs}.
\end{tutorial}

\section{Construction~: les estimateurs de covariance}\label{bld:qm:covariance-estimation-and-random-matrices:covest}

\begin{build}
\bfield{Rôle} Toutes les covariances utilisées par la construction de portefeuilles et la gestion du risque de la firme miniature (à partir du chapitre~23) proviennent de ce module, et son ratio entre volatilité réalisée et volatilité prévue est
suivi.
\bfield{Interface} \texttt{sample\_cov(X)}~; \texttt{mp\_edges(q)}, \texttt{mp\_density(x, q)}~; \texttt{lw\_identity(X)}, \texttt{lw\_constant\_corr(X)}~; \texttt{nonlinear\_shrinkage(X)}~;
\texttt{clip(X)}~; \texttt{pca\_factor(X, k)}~; \texttt{min\_var\_weights(Sigma)}~; \texttt{portfolio\_risk(w, Sigma)}.
\bfield{Règles} Ne jamais inverser une covariance empirique brute lorsque $N/T$ dépasse quelques pour cent~; indiquer l'intensité ou le seuil d'écrêtage avec chaque estimation~; juger les \omterm{def:qm:estimation:estimator}{estimateurs} sur
le risque hors échantillon des portefeuilles construits à partir d'eux, et non sur leur ajustement en échantillon.
\bfield{Tests d'acceptation} \texttt{code/firm/covest/tests/}~: la densité de Marchenko--Pastur a pour intégrale un et borne le spectre du bruit~; l'intensité de Ledoit--Wolf correspond à sa
définition par force brute et améliore l'erreur de Frobenius~; la cible à corrélation constante améliore un modèle vrai à corrélation constante~; le \omterm{def:qm:covariance-estimation-and-random-matrices:nonlinear}{rétrécissement non linéaire} suit les valeurs propres oracle~; l'écrêtage et le
\omterm{def:qm:covariance-estimation-and-random-matrices:factor}{modèle factoriel} conservent les variances empiriques~; les pondérations de variance minimale pour une covariance diagonale.
\bfield{Pour aller plus loin} POET (facteurs plus seuillage)~; \omterm{def:qm:covariance-estimation-and-random-matrices:nonlinear}{rétrécissement non linéaire} pour $N > T$~; un suivi glissant du ratio entre volatilité réalisée et volatilité prévue sur des portefeuilles en production.
\end{build}

\begin{omsources}
\item V.~A.~Marčenko et L.~A.~Pastur, ``Distribution of eigenvalues for some sets of random matrices'', \emph{Mathematics of the USSR-Sbornik} 1, 1967.
\item L.~Laloux, P.~Cizeau, J.-P.~Bouchaud et M.~Potters, ``Noise dressing of financial correlation matrices'', \emph{Physical Review Letters} 83, 1999.
\item O.~Ledoit et M.~Wolf, ``A well-conditioned estimator for large-dimensional covariance matrices'', \emph{Journal of Multivariate Analysis} 88, 2004; ``Honey, I shrunk the sample
covariance matrix'', \emph{Journal of Portfolio Management} 30, 2004.
\item O.~Ledoit et M.~Wolf, ``Analytical nonlinear shrinkage of large-dimensional covariance matrices'', \emph{Annals of Statistics} 48, 2020.
\item I.~M.~Johnstone, ``On the distribution of the largest eigenvalue in principal components analysis'', \emph{Annals of Statistics} 29, 2001.
\item J.~Baik, G.~Ben Arous et S.~Péché, \emph{Annals of Probability} 33, 2005~; J.~Baik et J.~W.~Silverstein, \emph{Journal of Multivariate Analysis} 97, 2006.
\item J.~Fan, Y.~Liao et M.~Mincheva, ``Large covariance estimation by thresholding principal orthogonal complements'', \emph{Journal of the Royal Statistical Society B} 75, 2013.
\item J.~Bun, J.-P.~Bouchaud et M.~Potters, ``Cleaning large correlation matrices: tools from random matrix theory'', \emph{Physics Reports} 666, 2017.
\end{omsources}

\section{Exercices}

\begin{exercise}[$\star$]\label{exo:qm:covariance-estimation-and-random-matrices:1}
Calculer les bords du spectre de Marchenko--Pastur pour $N = 500$ actions et $T = 1\,250$ jours~?
\end{exercise}

\begin{exercise}[$\star$]\label{exo:qm:covariance-estimation-and-random-matrices:2}
Un \omterm{def:qm:covariance-estimation-and-random-matrices:pca}{portefeuille de variance minimale} échantillonné, construit à partir de 300 actions et de 600 jours, annonce 5~\%. Quelle volatilité faut-il attendre~?
\end{exercise}

\begin{exercise}[$\star$]\label{exo:qm:covariance-estimation-and-random-matrices:3}
Montrer que la variance du \omterm{def:qm:covariance-estimation-and-random-matrices:pca}{portefeuille de variance minimale} vaut $1/(\mathbf 1^\top\Sigma^{-1}\mathbf 1)$.
\end{exercise}

\begin{exercise}[$\star\star$]\label{exo:qm:covariance-estimation-and-random-matrices:4}
Avec $q = 0.25$, quels pics de la population sont détectables, et où un pic de 3 apparaît-il dans l'échantillon~?
\end{exercise}

\begin{exercise}[$\star\star$]\label{exo:qm:covariance-estimation-and-random-matrices:5}
Montrer que le \omterm{def:qm:covariance-estimation-and-random-matrices:linear}{rétrécissement linéaire} $\delta mI + (1 - \delta)S$ conserve les vecteurs propres de $S$ et envoie chaque valeur propre $\lambda$ sur $\delta m + (1 - \delta)\lambda$. Pourquoi est-il moins souple que le
\omterm{def:qm:covariance-estimation-and-random-matrices:nonlinear}{rétrécissement non linéaire}~?
\end{exercise}

\begin{exercise}[$\star\star$]\label{exo:qm:covariance-estimation-and-random-matrices:6}
Combien de jours de données faudrait-il pour que le \omterm{def:qm:covariance-estimation-and-random-matrices:pca}{portefeuille de variance minimale} échantillonné de 200 actions sous-estime sa volatilité de 10~\% au plus~?
\end{exercise}

\begin{exercise}[$\star\star\star$]\label{exo:qm:covariance-estimation-and-random-matrices:7}
\emph{Programmation.} Reproduire le rapport entre la volatilité réelle et la volatilité prévue du \omterm{def:qm:covariance-estimation-and-random-matrices:pca}{portefeuille de variance minimale} échantillonné pour $q$ allant de 0.1 à 0.9, et le comparer à $1/(1 - q)$.
\end{exercise}

\begin{exercise}[$\star\star\star$]\label{exo:qm:covariance-estimation-and-random-matrices:8}
\emph{Trouver l'erreur.} «~Nous avons choisi notre \omterm{def:qm:estimation:estimator}{estimateur} de covariance parce qu'il présente la plus faible volatilité, dans l'échantillon, du \omterm{def:qm:covariance-estimation-and-random-matrices:pca}{portefeuille de variance minimale} parmi tous les \omterm{def:qm:estimation:estimator}{estimateurs} que nous avons essayés.~»
\end{exercise}

\section{Problème~: le portefeuille qui promettait 6~\%}

\begin{problem}[{Problème du week-end --- risque promis et risque livré}]\label{pb:qm:covariance-estimation-and-random-matrices:1}
Deux cents actions suivent une covariance connue, avec un facteur de marché et cinq secteurs de quarante. Un gérant estime la covariance à partir de deux ans (500 jours) de rendements quotidiens, construit le
\omterm{def:qm:covariance-estimation-and-random-matrices:pca}{portefeuille de variance minimale} et le conserve pendant un an. Cinquante historiques sont simulés.

\medskip\noindent\textbf{Partie I --- L'échantillon.}
\begin{enumerate}
\item Que prévoit et que livre le portefeuille échantillonné, et quel est le véritable optimum~?
\item Quel rapport la formule $1/(1 - q)$ prévoit-elle, et que donne la simulation~?
\item Quels sont les bords de Marchenko--Pastur pour $q = 0.4$, et où tombent les valeurs propres de pur bruit~?
\item Quelles valeurs propres du \omterm{def:qm:covariance-estimation-and-random-matrices:factor}{modèle factoriel} se détachent, et de combien sont-elles biaisées~?
\item Quel est le plus petit facteur détectable à ce $q$~?
\end{enumerate}

\medskip\noindent\textbf{Partie II --- Meilleurs \omterm{def:qm:estimation:estimator}{estimateurs}.}
\begin{enumerate}[resume]
\item Quelles intensités les deux \omterm{def:qm:covariance-estimation-and-random-matrices:linear}{rétrécissements linéaires} choisissent-ils, et que livrent-ils~?
\item Que livrent le \omterm{def:qm:covariance-estimation-and-random-matrices:nonlinear}{rétrécissement non linéaire} et l'écrêtage~?
\item Que livre le modèle à six facteurs, et pourquoi est-il le meilleur ici~?
\item Quels sont les rapports réalisé sur prévu des six~?
\item Quelle prévision d'\omterm{def:qm:estimation:estimator}{estimateur} est la plus honnête, et quel portefeuille est le meilleur~?
\end{enumerate}

\medskip\noindent\textbf{Partie III --- Dimensions.}
\begin{enumerate}[resume]
\item Que devient le portefeuille échantillonné à $q = 0.9$~?
\item Combien de jours faudrait-il pour que le rapport vaille 1.1~?
\item Pourquoi l'optimiseur choisit-il les directions les plus bruitées~?
\item Que fait une contrainte de positions longues uniquement à ce problème~?
\item Pourquoi l'identité constitue-t-elle une mauvaise cible de rétrécissement pour les actions~?
\end{enumerate}

\medskip\noindent\textbf{Partie IV --- Jugement.}
\begin{enumerate}[resume]
\item Quel \omterm{def:qm:estimation:estimator}{estimateur} faudrait-il mettre en production, et pourquoi~?
\item Comment faudrait-il le surveiller~?
\item Qu'est-ce qui changerait avec des rendements à queues épaisses ou une covariance variable~?
\item Énoncer le \emph{résultat nommé}~: le rapport risque réalisé sur risque prévu du portefeuille échantillonné à $q = 0.4$, comparé à $1/(1 - q)$, et le rapport atteint par chacun des meilleurs \omterm{def:qm:estimation:estimator}{estimateurs}.
\item En une phrase~: pourquoi un optimiseur aggrave-t-il l'erreur d'estimation~?
\end{enumerate}
\end{problem}

\section{Questions d'entretien}

\begin{interviewq}[$\star$ \iqroles{chercheur, risque}]\label{iq:qm:covariance-estimation-and-random-matrices:1}
Pourquoi un \omterm{def:qm:covariance-estimation-and-random-matrices:pca}{portefeuille de variance minimale} construit à partir d'une covariance échantillonnée sous-estime-t-il son risque~?
\end{interviewq}

\begin{interviewq}[$\star\star$ \iqroles{chercheur, ingénieur ML}]\label{iq:qm:covariance-estimation-and-random-matrices:2}
Qu'est-ce que la \omterm{thm:qm:covariance-estimation-and-random-matrices:mp}{loi de Marchenko--Pastur}, et comment l'utilise-t-on sur une matrice de corrélation~?
\end{interviewq}

\begin{interviewq}[$\star\star$ \iqroles{chercheur}]\label{iq:qm:covariance-estimation-and-random-matrices:3}
Expliquer le rétrécissement de Ledoit--Wolf. De quoi dépend l'intensité~?
\end{interviewq}

\begin{interviewq}[$\star\star$ \iqroles{chercheur, risque}]\label{iq:qm:covariance-estimation-and-random-matrices:4}
Vous disposez de 1\,000 actions et de trois ans de données quotidiennes. Comment estimez-vous la \omterm{def:qm:covariance-estimation-and-random-matrices:cov}{matrice de covariance}~?
\end{interviewq}

\begin{interviewq}[$\star\star$ \iqroles{ingénieur ML, chercheur}]\label{iq:qm:covariance-estimation-and-random-matrices:5}
Combien de composantes principales d'une matrice de corrélation d'actions portent de l'information~?
\end{interviewq}

\begin{interviewq}[$\star\star\star$ \iqroles{chercheur}]\label{iq:qm:covariance-estimation-and-random-matrices:6}
Qu'est-ce qu'un \omterm{def:qm:covariance-estimation-and-random-matrices:nonlinear}{estimateur équivariant par rotation}, et que peut-on faire de mieux dans cette classe~?
\end{interviewq}
